% Chapter 6: Architecture
% The chapter's own opening text goes after \chapter; its sections are files in
% chapters/6-architecture/, added below with \input in the order they appear in the book.

\chapter{Architecture}

% Draft by Claude, 2026-10-24, from webdoc internals pages (index, storage, life-of-a-step) and the
% server code; the queue-log bytes below are from a 2.6.0 node.
A Queen node is one process. Every write is decided once, on the leader, and then applied on every
node. This chapter follows a write through the process, from the request to the disk.

The rule behind the whole design is: \textbf{one leader per tenant}. A tenant lives in one raft
group, and one leader orders all its writes. So a transaction, even across ten queues, is one log
entry: no coordinator, no two-phase commit. The price is that one tenant's writes do not spread
over more nodes. To use more nodes for writes, run more raft groups (see Many tenants on one
cluster): with as many groups as nodes, every node leads one.

\queenfig{raft-groups}{One leader per tenant: three raft groups, each led on a different node.}

% Chapter 6, section 1: Leader and followers
% Write in plain LaTeX. Code of a guide step: \guide{<step>}; commands: \begin{shell}...\end{shell};
% a screenshot: \screenshot{file}{caption}; a gap in the draft: \placeholder{...}.
% See book/latex/queen.sty for the full list.

\section{Leader and followers}

A raft group has one leader. The other nodes are followers, and every node has a full copy of the
data.

The leader orders all the writes. A follower is not idle: it reads the requests of its own
clients, prepares each one, and sends the leader only the prepared command. When the write is
committed, the follower answers its client from its own copy. So the leader spends its time
ordering and writing, not parsing JSON.

\queenfig{follower-offload}{A follower does the work of its own clients.}

A client can talk to any node. A read waits until that node has applied everything that was
committed when the read started, so nobody reads behind a write they were already told about.

If the leader stops with SIGTERM, it hands its role to the most up-to-date follower, in a fraction
of a second. If it crashes, the others wait one to two seconds and elect a new leader.

% Chapter 6, section 2: Admission, Batch, Plan, Apply
% Write in plain LaTeX. Code of a guide step: \guide{<step>}; commands: \begin{shell}...\end{shell};
% a screenshot: \screenshot{file}{caption}; a gap in the draft: \placeholder{...}.
% See book/latex/queen.sty for the full list.

\section{Admission, Batch, Plan, Apply}

A write goes through four steps.

\queenfig{write-path}{The four steps of a write.}

\paragraph{Admission.} Every node has a memory budget for the writes waiting to be planned (128 MiB
by default). A push takes its share before its body is read. When the budget is full, the push
waits, so a fast producer just slows down. After about 15 seconds it gets a 429, and the client
tries again. Pops and acks never wait here. We added this after a broker filled 10 GB of memory
with waiting requests and was killed.

\paragraph{Batch.} The leader collects up to 4,096 commands (or 4 MiB) and plans them together, as
one log entry. Up to eight entries are in flight at the same time. A retried request is found by
its id and gets the answer it missed, so nothing is planned twice.

\paragraph{Plan.} The planner decides what each command means right now, and writes the answer
down as facts: ``append these messages to partition 17, at offsets 4,210 to 4,309, at this
microsecond''. Ids, offsets and times are chosen here, once. To use more cores, eight lanes plan in
parallel: each lane owns a share of the partitions, and one last step plans what touches more than
one lane.

\queenfig{lanes}{Eight lanes plan in parallel; one step plans what crosses lanes.}

\paragraph{Apply.} Every node executes the committed entries, in order. Apply reads no clock and no
random number, so every node ends with the same bytes. Four threads apply the work of different
partitions at the same time.

% Chapter 6, section 3: Raft internal protocol
% Write in plain LaTeX. Code of a guide step: \guide{<step>}; commands: \begin{shell}...\end{shell};
% a screenshot: \screenshot{file}{caption}; a gap in the draft: \placeholder{...}.
% See book/latex/queen.sty for the full list.

\section{Raft internal protocol}

The nodes agree on the log with Raft, through the openraft library. Raft decides who leads and when
an entry is committed. An entry is committed when a majority of the nodes has written it and
synced it to disk (fsync).

Queen has no separate Raft log: the queue log files are the log (see Storage). The nodes talk on
port 7400, over plain HTTP, with a shared token. The leader sends a heartbeat every 100~ms. If the
followers hear nothing for one to two seconds, they elect a new leader.

A leader that hears from no majority for 1.5 seconds stops logging new writes and keeps them
waiting, in order. After 3 seconds it hands leadership to another node.

\queenfig{raft-timings}{The timings Raft runs on.}

A node that was away gets the entries it missed from the leader, who keeps them for 6 hours while
its disk has room. After that, the leader sends a snapshot instead: since 2.6, only the part the
node is missing. The node exits with code 75 and loads the snapshot when it starts again.

% Chapter 6, section 4: Memory state
% Write in plain LaTeX. Code of a guide step: \guide{<step>}; commands: \begin{shell}...\end{shell};
% a screenshot: \screenshot{file}{caption}; a gap in the draft: \placeholder{...}.
% See book/latex/queen.sty for the full list.

\section{Memory state}

Each node keeps its whole state in memory: queues, partitions, consumer positions, KV, timers,
request ids, and the dedup records of recent pushes. It is loaded once, at boot.

Messages are not in memory. They stay in the queue log files and are read from there. So memory
grows with the number of partitions and keys, and with the pushes of the dedup window, not with how
many messages a queue keeps.

\queenfig{memory-disk}{What lives in memory, and what lives on disk.}

Message traces are kept apart, in their own file, and are read from disk.

% Chapter 6, section 5: Consumption engine
% Write in plain LaTeX. Code of a guide step: \guide{<step>}; commands: \begin{shell}...\end{shell};
% a screenshot: \screenshot{file}{caption}; a gap in the draft: \placeholder{...}.
% See book/latex/queen.sty for the full list.

\section{Consumption engine}

Pops, acks and leases do not go through the planner. The leader keeps each consumer group's
positions and leases in memory, and gives out the oldest ready partition first.

The positions that changed go into the log as one small checkpoint, at once when an answer is
waiting and at least every 5~ms. A pop gets its answer only after that checkpoint is committed. A
pop that finds nothing waits, and is woken when a new message arrives in its partition.

\queenfig{consumption-engine}{Pops, acks and leases, served from the leader's memory.}

We moved consumption out of the planner because acks were waiting behind pushes: at 300,000
messages a second, an ack took more than 500~ms. Now many acks travel in one checkpoint.

A partition idle for 10 minutes leaves memory, and is loaded again with its next message. A new
leader waits about half a second before it serves pops, so the same messages are never given out
twice.

% Chapter 6, section 6: LMDB
% Write in plain LaTeX. Code of a guide step: \guide{<step>}; commands: \begin{shell}...\end{shell};
% a screenshot: \screenshot{file}{caption}; a gap in the draft: \placeholder{...}.
% See book/latex/queen.sty for the full list.

\section{LMDB}

The state in memory is saved to an LMDB file, \texttt{store/data.mdb}, four times a second. This is
a checkpoint, and no client waits for it: it only decides how much of the log a restart has to
replay.

Every value carries a checksum, and the node checks all of them at boot. Every 6 hours a
background scan reads the whole store again, to find damage before a restart does.

Traces have their own LMDB file, \texttt{traces/data.mdb}.

% Chapter 6, section 7: Storage
% Write in plain LaTeX. Code of a guide step: \guide{<step>}; commands: \begin{shell}...\end{shell};
% a screenshot: \screenshot{file}{caption}; a gap in the draft: \placeholder{...}.
% See book/latex/queen.sty for the full list.

\section{Storage}

A message is written once, in its queue's log files, and that write is what makes it durable:
there is no second copy. We built it this way on purpose. An earlier design wrote every message
twice, first in a consensus log and then in a segment file. Now the bytes that make a message
durable are the bytes a pop reads back.

\subsection{The data directory}

Everything a node knows is in one directory:

\begin{shell}
/var/lib/queen/raft
  qlog/                    the messages: one log per queue
    LANES                  the log layout, fixed at creation (and SHARDS, when set)
    q0/                    the system log: entries that touch no queue
    q5837314263283300749/  one directory per queue
      r00000001.qlog       a log file, up to 64 MiB, then the next one starts
      r00000001.qidx       the index of a closed file
      r00000001.qhsh       the dedup hashes of a closed file, for a fast restart
  store/data.mdb           the state checkpoint (LMDB)
  traces/data.mdb          message traces (LMDB)
  raft/                    cluster only: vote, members, what is kept for an away node
  snapshots/               a snapshot from the leader, waiting for the next start
  seg/  log/               almost empty, from an earlier design
  local.db  dash.db        this node's own metrics and dashboard data
  groups/g1/               the same tree for each extra raft group
\end{shell}

A closed file never changes again. Payloads of 512 bytes or more are compressed with zstd, when it
saves at least a tenth. With thousands of queues, \texttt{QUEEN\_QLOG\_SHARDS} puts them all into a
few shared logs, so a write syncs a few files instead of thousands.

\subsection{What is inside a log file}

A log file is a 32-byte header and then records, one after the other. Here is the start of a real
file, after one push of \texttt{order-8891} with data \texttt{\{"orderId":8891\}}:

\begin{shell}
xxd qlog/q5837314263283300749/r00000001.qlog | head -10
00000000: 514e 514c 4f47 3100 0100 0000 0000 0000  QNQLOG1.........
00000010: 0b00 0000 0000 0000 0000 0000 0000 0000  ................
00000020: 6e00 0000 da25 b6dc 1ed0 878c 0b00 0000  n....%..........
00000030: 0000 0000 0100 0000 0000 0000 0000 0000  ................
00000040: 0000 0000 0100 0000 75e7 b953 935e 0600  ........u..S.^..
00000050: 00b9 7eb7 7dcc 2db6 dcbc ff49 4354 80b3  ..~.}.-....ICT..
00000060: c82d 0000 0000 01a1 6d65 4e4f 7001 bf27  .-......meNOp..'
00000070: b56a 9334 d445 0a00 6f72 6465 722d 3838  .j.4.E..order-88
00000080: 3931 7b22 6f72 6465 7249 6422 3a38 3839  91{"orderId":889
00000090: 317d da01 0000 6b08 25f1 17ae 81e2 0b00  1}....k.%.......
\end{shell}

The record starts at byte 0x20. Every number is little-endian.

\begin{options}{Field}{Bytes}{In this record}
len & 4 & 110: the bytes after this field. A reader always knows where the next record starts. \\
checksum & 8 & An xxh3 of everything after it. A record torn by a crash fails it, and is cut. \\
seq & 8 & 11: the log entry that wrote this record. \\
partition & 8 & 1: the partition id. \\
first offset & 8 & 0. \\
count & 4 & 1 message. \\
created at & 8 & The time of the push, in microseconds. \\
kind & 1 & 0: messages. The top bit says the payload is compressed. \\
hashes & 16 a message & The hash of \texttt{order-8891}, for dedup. \\
payload & the rest & One frame per message: its length (45), flags, its 16-byte id, its transactionId and its data. \\
\end{options}

\queenfig{record-layout}{The message record of the dump, field by field.}

The next record (474 bytes, at 0x92) has the same seq, 11: it is the log entry itself, without the
payload. An entry that touches several queues is written into each of their logs, and after a
crash it is replayed only if every copy is there. This is how a transaction stays all or nothing on
disk.

\queenfig{entry-copies}{An entry in every log it touches; replayed only if every copy is there.}

\subsection{How retention works}

Retention is decided once, on the leader, like every write. A scanner walks all the partitions
(50,000 a second) and proposes what can go. The planner checks each proposal and logs watermarks:
below one the messages are gone, below another their dedup records are gone. A partition that is
empty, with no dead letter, no lease and no write for the cleanup age, is deleted.

Moving a watermark frees no disk by itself. Each node then frees its own files, one whole file at a
time.

\queenfig{retention-watermarks}{Retention decided on the leader, as watermarks.}

\subsection{Rewriting a file}

A closed file with no live record is deleted. A closed file that is partly live is rewritten: the
node copies only the live records into a new file, without stopping the writes, and swaps it in.
A queue's own file is rewritten at its first dead message, a shared file when half of it is dead.

The active file is never touched, and neither is a file that recovery or a node that is behind
still needs. One pass looks at up to 256 files, for at most 200~ms, and rewrites at most one file;
while work is left, the next pass is four times bigger. An active file with data older than 10
minutes is closed early, so retention can reach it, when it holds a dead message or is an eighth
full.

\queenfig{file-rewrite}{Reclaim: dead files deleted, partly live files rewritten.}

\why{a file is the unit the disk gives back. Rewriting a file with only its live records gives back
the space of the dead ones without waiting for the whole file to die, and the small passes keep the
rewrite from slowing down the writes.}

